\section{Mathematical Definition and Properties}
\label{sec:definition}

\subsection{The Uniform Swing Seats-Votes Curve}

For a $k$-district plan with observed Democratic vote shares $v_1, \ldots, v_k$
(from a baseline election) and observed statewide mean $\bar{v}_D = k^{-1}\sum_i v_i$,
the uniform swing seats-votes curve is:
\[
S(v) = \frac{1}{k}\sum_{i=1}^k \mathbf{1}[v_i + (v - \bar{v}_D) > 0.50],
\quad v \in [0,1].
\]
$S(v)$ is a step function taking values in $\{0, 1/k, 2/k, \ldots, 1\}$.
It increases by $1/k$ at each value $v^* = 0.50 - v_i + \bar{v}_D$ where
a district would flip from Republican to Democratic.

The curve is well-defined for all $v$ if no district has exactly $v_i = \bar{v}_D
- (v - 0.50)$, which occurs with probability zero for continuous vote shares.
In practice, district vote shares are discrete (counts), so ties are broken
using the district rank in the sorted vote-share sequence.

\subsection{Responsiveness}

Responsiveness is the slope of $S(v)$ at the 50\% vote-share pivot:
\[
R = \left.\frac{dS}{dv}\right|_{v=0.50} = \frac{1}{k \cdot \Delta v_{\text{pivot}}},
\]
where $\Delta v_{\text{pivot}}$ is the vote-share interval around 50\% within which
the next district flips. For a plan with $m$ districts near the pivot,
$R \approx m / (k \cdot \Delta v)$. Intuitively, $R$ counts how many districts
would flip per unit of vote-share change near the pivot.

A plan with $R = 2.0$ has the property that a 1-point statewide vote swing
near 50\% changes the seat share by 2 points---the competitive ideal. This
translates to: for NC ($k=14$), a 1-point vote swing flips approximately
$14 \times 2 / 100 \approx 0.28$ districts, i.e., one district for every
$\approx 3.5$-point swing. Plans with $R < 1.5$ have fewer competitive
districts, so the same swing flips fewer seats, insulating incumbents from
electoral accountability.

\subsection{The Sufficiency Result}

The seats-votes curve is the sufficient statistic for the two most-used scalar
partisan metrics: Efficiency Gap and Partisan Bias.

\textbf{EG from $S(v)$.} The Efficiency Gap, under the uniform distribution
of statewide vote shares, equals the signed area between the $S(v)$ curve
and the diagonal $S(v) = v$:
\[
\text{EG} = \int_0^1 [S(v) - v]\, dv \approx \bar{S} - \bar{v}_D,
\]
where the approximation holds when the uniform distribution over $[0,1]$ is
a reasonable model for the distribution of statewide vote shares. More precisely,
$\text{EG} = s - 2v + 1/2$ (the vote-share formulation from L.1), which equals
the signed area under the exact step function $S(v)$.

\textbf{Bias from $S(v)$.} Partisan Bias is simply $S(0.50) - 0.50$---the
value of the $S(v)$ curve at the 50\% pivot, minus 0.50. Reading Partisan Bias
from the seats-votes curve requires only locating the point $(0.50, S(0.50))$
and measuring its vertical distance from $(0.50, 0.50)$.

\textbf{Corollary.} A court that admits the seats-votes curve as evidence
already has access to both EG and Partisan Bias as scalar summaries of that
curve. The reverse is not true: EG and Partisan Bias as scalars do not
uniquely determine the seats-votes curve. The curve is therefore strictly
more informative than any combination of its scalar summaries.

\subsection{The Competitive Ideal}

The symmetric seats-votes curve---$S(v) = v$ for all $v$, the diagonal---would
imply perfect proportionality: a party winning $v\%$ of votes wins $v\%$ of
seats. This is not a realistic target for single-member district systems, which
structurally produce a majority-bonus ($S$ above the diagonal for vote shares
above 50\% and below the diagonal for vote shares below 50\%). The competitive
ideal is not the diagonal but a curve passing through $(0.50, 0.50)$ with slope
$R \approx 2.0$---the curve that translates vote swings into seat changes at
a competitive rate without concentrating seat changes in a narrow vote-share band
that insulates incumbents.
